\documentclass[12pt,a4paper]{article}

% ------------------------------------------------
% PAQUETES
% ------------------------------------------------
\usepackage[utf8]{inputenc}
\usepackage[T1]{fontenc}
\usepackage[spanish]{babel}

\usepackage{geometry}
\geometry{margin=2.5cm}

\usepackage{longtable}
\usepackage{array}
\usepackage{booktabs}
\usepackage{tabularx}
\usepackage{titlesec}
\usepackage{graphicx}
\usepackage{float}
\usepackage{hyperref}
\usepackage[table]{xcolor}
\usepackage{fancyhdr}

\definecolor{brandyellow}{RGB}{245,225,96}
\definecolor{lightyellow}{RGB}{255,249,196}
\definecolor{darkgray}{RGB}{40,40,40}
\definecolor{linkblue}{RGB}{0,76,153}

\hypersetup{
    colorlinks=true,
    linkcolor=linkblue,
    urlcolor=linkblue,
    citecolor=linkblue
}

% ------------------------------------------------
% CONFIGURACIÓN
% ------------------------------------------------
\setlength{\parskip}{0.5em}
\setlength{\parindent}{0pt}
\setlength{\tabcolsep}{8pt}
\arrayrulecolor{darkgray}

\titleformat{\section}{\normalfont\Large\bfseries\color{black}}{\thesection}{1em}{}
\titlespacing*{\section}{0pt}{*3}{*1.2}
\titleformat{\subsection}{\normalfont\large\bfseries\color{darkgray}}{\thesubsection}{1em}{}
\titlespacing*{\subsection}{0pt}{*2}{*1}

\title{Documentación de Base de Datos \\ \large RecipeBook}
\author{
    Giovanny Sierra Reina \\
    Angel Andres Díaz Vergara \\
    Juan Felipe Guevara Olaya \\
    Mathew Stev Toro Mondragón \\
    David Sánchez Acero
}
\date{27/05/2026}

% ------------------------------------------------
% DOCUMENTO
% ------------------------------------------------
\begin{document}

\pagestyle{fancy}
\fancyhf{}
\renewcommand{\headrulewidth}{0pt}
\renewcommand{\footrulewidth}{0pt}
\fancyhead[L]{\textcolor{darkgray}{RecipeBook}}
\fancyhead[R]{\textcolor{darkgray}{Documentación de Base de Datos}}
\fancyfoot[C]{\textcolor{darkgray}{\thepage}}
\setlength{\headheight}{14pt}
\rowcolors{2}{lightyellow!60}{white}

% ------------------------------------------------
% PORTADA
% ------------------------------------------------
\begin{titlepage}
    \centering

    \vspace*{2cm}

    \colorbox{lightyellow}{\parbox{0.9\textwidth}{\centering
        {\Huge\bfseries\textcolor{black}{RecipeBook}}\\[0.5em]
        {\Large\textcolor{darkgray}{Documentación de Base de Datos}}\\[0.5em]
    }}

    \vspace{1.5cm}

    \vspace{2cm}

    \textbf{Integrantes del equipo:}

    \vspace{0.5cm}

    Giovanny Sierra Reina\\
    Angel Andres Díaz Vergara\\
    Juan Felipe Guevara Olaya\\
    Mathew Stev Toro Mondragón\\
    David Sánchez Acero

    \vfill

    {\Large Universidad Distrital Francisco José de Caldas\par}

    \vspace{0.3cm}

    Bases de Datos\\
    2026

\end{titlepage}

% ------------------------------------------------
% TABLA DE VERSIONAMIENTO
% ------------------------------------------------
\section*{Control de Versiones}

\begin{table}[H]
\centering
\begin{tabular}{|c|c|c|p{6cm}|}
\hline
\textbf{Versión} & \textbf{Fecha} & \textbf{Autor} & \textbf{Descripción del cambio} \\
\hline
1.0 & 27/05/2026 & Equipo RecipeBook & Creación inicial del documento \\
\hline
1.1 & --/--/---- & --- & --- \\
\hline
\end{tabular}
\end{table}

\vspace{1cm}

% ------------------------------------------------
% ÍNDICE
% ------------------------------------------------
\tableofcontents
\newpage

% ------------------------------------------------
% DESCRIPCIÓN DEL SISTEMA
% ------------------------------------------------
\section{Descripción del Sistema de Información}

RecipeBook es una aplicación web orientada a la gestión de recetas culinarias. El sistema permite registrar, consultar y administrar información relacionada con recetas, ingredientes, usuarios y categorías.

La plataforma fue desarrollada utilizando Java, JSP, Servlets y PostgreSQL, implementando una arquitectura DAO mediante JDBC para el acceso a datos.

% ------------------------------------------------
% DESCRIPCIÓN DE LA BASE DE DATOS
% ------------------------------------------------
\section{Descripción de la Base de Datos}

La base de datos del proyecto fue diseñada bajo un modelo relacional con el objetivo de garantizar integridad, consistencia y organización de la información.

Características principales:

\begin{itemize}
    \item Uso de claves primarias y foráneas para mantener la integridad referencial.
    \item Restricciones de integridad como \texttt{NOT NULL}, \texttt{UNIQUE} y \texttt{CHECK}.
    \item Relaciones bien definidas entre entidades de usuario, recetas, pasos, ingredientes y categorías.
    \item Persistencia gestionada en PostgreSQL.
    \item Acceso a datos mediante JDBC directo y un patrón DAO centralizado en \texttt{DAOFactory}.
\end{itemize}

% ------------------------------------------------
% DICCIONARIO DE DATOS
% ------------------------------------------------
\section{Diccionario de Datos}

\subsection{Tabla: TIPO\_RECETA}

Esta tabla almacena los diferentes tipos de recetas disponibles en la plataforma.

\renewcommand{\arraystretch}{1.3}
\small
\begin{longtable}{|p{3.5cm}|p{2.5cm}|p{3.5cm}|p{5cm}|}
\hline
\textbf{Campo} & \textbf{Tipo de dato} & \textbf{Restricciones} & \textbf{Propósito / Descripción} \\
\hline
id\_tipo\_receta & SERIAL & PRIMARY KEY, NOT NULL & Identificador único del tipo de receta. \\
\hline
nombre\_tipo & VARCHAR(50) & NOT NULL, UNIQUE & Nombre del tipo de receta (por ejemplo: desayuno, almuerzo). \\
\hline
descripcion & TEXT & NULL & Descripción opcional del tipo de receta. \\
\hline
\end{longtable}

\subsection{Tabla: MULTIMEDIA}

Esta tabla centraliza referencias a recursos multimedia asociados a usuarios, recetas y pasos.

\begin{longtable}{|p{3.5cm}|p{2.5cm}|p{3.5cm}|p{5cm}|}
\hline
\textbf{Campo} & \textbf{Tipo de dato} & \textbf{Restricciones} & \textbf{Propósito / Descripción} \\
\hline
id\_multimedia & SERIAL & PRIMARY KEY, NOT NULL & Identificador único del recurso multimedia. \\
\hline
url & VARCHAR(255) & NOT NULL, UNIQUE & Ruta o URL del recurso multimedia. \\
\hline
\end{longtable}

\subsection{Tabla: ROL}

Define los roles de usuario que se usan para control de permisos.

\begin{longtable}{|p{3.5cm}|p{2.5cm}|p{3.5cm}|p{5cm}|}
\hline
\textbf{Campo} & \textbf{Tipo de dato} & \textbf{Restricciones} & \textbf{Propósito / Descripción} \\
\hline
id\_rol & SERIAL & PRIMARY KEY, NOT NULL & Identificador único del rol. \\
\hline
nombre\_rol & VARCHAR(50) & NOT NULL, UNIQUE & Nombre del rol asignado al usuario. \\
\hline
\end{longtable}

\subsection{Tabla: MODULO}

Contiene los módulos del sistema asociados a operaciones de seguridad o navegación.

\begin{longtable}{|p{3.5cm}|p{2.5cm}|p{3.5cm}|p{5cm}|}
\hline
\textbf{Campo} & \textbf{Tipo de dato} & \textbf{Restricciones} & \textbf{Propósito / Descripción} \\
\hline
id\_modulo & SERIAL & PRIMARY KEY, NOT NULL & Identificador del módulo. \\
\hline
nombre\_modulo & VARCHAR(50) & NOT NULL, UNIQUE & Nombre del módulo del sistema. \\
\hline
\end{longtable}

\subsection{Tabla: OPERACION}

Relaciona acciones del sistema con módulos específicos.

\begin{longtable}{|p{3.5cm}|p{2.5cm}|p{3.5cm}|p{5cm}|}
\hline
\textbf{Campo} & \textbf{Tipo de dato} & \textbf{Restricciones} & \textbf{Propósito / Descripción} \\
\hline
id\_operacion & SERIAL & PRIMARY KEY, NOT NULL & Identificador de la operación. \\
\hline
accion & VARCHAR(100) & NOT NULL, UNIQUE & Nombre de la acción o permiso. \\
\hline
id\_modulo & INT & NOT NULL, FOREIGN KEY & Referencia al módulo al que pertenece la operación. \\
\hline
\end{longtable}

\subsection{Tabla: ROL\_OPERACION}

Tabla de relación muchos a muchos entre roles y operaciones.

\begin{longtable}{|p{3.5cm}|p{2.5cm}|p{3.5cm}|p{5cm}|}
\hline
\textbf{Campo} & \textbf{Tipo de dato} & \textbf{Restricciones} & \textbf{Propósito / Descripción} \\
\hline
id\_rol & INT & NOT NULL, FOREIGN KEY & Referencia al rol. \\
\hline
id\_operacion & INT & NOT NULL, FOREIGN KEY & Referencia a la operación. \\
\hline
\end{longtable}

\subsection{Tabla: USUARIO}

Almacena los usuarios registrados en el sistema.

\begin{longtable}{|p{3.5cm}|p{2.5cm}|p{3.5cm}|p{5cm}|}
\hline
\textbf{Campo} & \textbf{Tipo de dato} & \textbf{Restricciones} & \textbf{Propósito / Descripción} \\
\hline
id\_usuario & SERIAL & PRIMARY KEY, NOT NULL & Identificador único del usuario. \\
\hline
nombre\_1 & VARCHAR(20) & NOT NULL & Primer nombre del usuario. \\
\hline
nombre\_2 & VARCHAR(20) & NULL & Segundo nombre del usuario (opcional). \\
\hline
apellido\_1 & VARCHAR(20) & NOT NULL & Primer apellido del usuario. \\
\hline
apellido\_2 & VARCHAR(20) & NULL & Segundo apellido del usuario (opcional). \\
\hline
correo\_electronico & VARCHAR(70) & NOT NULL, UNIQUE & Correo electrónico del usuario. \\
\hline
username & VARCHAR(100) & NOT NULL, UNIQUE & Nombre de usuario para el inicio de sesión. \\
\hline
password & VARCHAR(100) & NOT NULL & Contraseña cifrada o almacenada para autenticación. \\
\hline
id\_multimedia & INT & NULL, FOREIGN KEY & Referencia al recurso multimedia asociado al usuario. \\
\hline
id\_rol & INT & NOT NULL, FOREIGN KEY & Referencia al rol asignado al usuario. \\
\hline
\end{longtable}

\subsection{Tabla: RECETAS}

Contiene las recetas creadas por usuarios.

\begin{longtable}{|p{3.5cm}|p{2.5cm}|p{3.5cm}|p{5cm}|}
\hline
\textbf{Campo} & \textbf{Tipo de dato} & \textbf{Restricciones} & \textbf{Propósito / Descripción} \\
\hline
id\_receta & SERIAL & PRIMARY KEY, NOT NULL & Identificador único de la receta. \\
\hline
nombre\_receta & VARCHAR(100) & NOT NULL & Nombre de la receta. \\
\hline
descripcion & TEXT & NULL & Descripción detallada de la receta. \\
\hline
tiempo\_preparacion & INT & NULL & Tiempo estimado de preparación en minutos. \\
\hline
id\_usuario & INT & NOT NULL, FOREIGN KEY & Referencia al usuario autor de la receta. \\
\hline
id\_imagen & INT & NULL, FOREIGN KEY & Referencia a la imagen asociada a la receta. \\
\hline
\end{longtable}

\subsection{Tabla: PASOS}

Define los pasos de preparación para cada receta.

\begin{longtable}{|p{3.5cm}|p{2.5cm}|p{3.5cm}|p{5cm}|}
\hline
\textbf{Campo} & \textbf{Tipo de dato} & \textbf{Restricciones} & \textbf{Propósito / Descripción} \\
\hline
id\_paso & SERIAL & PRIMARY KEY, NOT NULL & Identificador único del paso. \\
\hline
id\_receta & INT & NOT NULL, FOREIGN KEY & Referencia a la receta asociada. \\
\hline
descripcion & TEXT & NOT NULL & Descripción del paso de preparación. \\
\hline
id\_multimedia & INT & NULL, FOREIGN KEY & Referencia a un recurso multimedia del paso. \\
\hline
\end{longtable}

\subsection{Tabla: INGREDIENTES}

Lista los ingredientes utilizados en cada receta o paso.

\begin{longtable}{|p{3.5cm}|p{2.5cm}|p{3.5cm}|p{5cm}|}
\hline
\textbf{Campo} & \textbf{Tipo de dato} & \textbf{Restricciones} & \textbf{Propósito / Descripción} \\
\hline
id\_ingrediente & SERIAL & PRIMARY KEY, NOT NULL & Identificador único del ingrediente. \\
\hline
id\_receta & INT & NOT NULL, FOREIGN KEY & Referencia a la receta donde se usa el ingrediente. \\
\hline
id\_paso & INT & NULL, FOREIGN KEY & Referencia al paso asociado (opcional). \\
\hline
nombre\_ingrediente & VARCHAR(255) & NOT NULL & Nombre del ingrediente. \\
\hline
\end{longtable}

\subsection{Tabla: UTENSILIOS}

Registra los utensilios necesarios para la preparación de cada receta.

\begin{longtable}{|p{3.5cm}|p{2.5cm}|p{3.5cm}|p{5cm}|}
\hline
\textbf{Campo} & \textbf{Tipo de dato} & \textbf{Restricciones} & \textbf{Propósito / Descripción} \\
\hline
id\_utensilio & SERIAL & PRIMARY KEY, NOT NULL & Identificador único del utensilio. \\
\hline
id\_receta & INT & NOT NULL, FOREIGN KEY & Referencia a la receta que requiere el utensilio. \\
\hline
id\_paso & INT & NOT NULL, FOREIGN KEY & Referencia al paso donde se usa el utensilio. \\
\hline
nombre\_utensilio & VARCHAR(255) & NOT NULL & Nombre del utensilio. \\
\hline
\end{longtable}

\subsection{Tabla: VALORACION}

Almacena las valoraciones que los usuarios hacen a las recetas.

\begin{longtable}{|p{3.5cm}|p{2.5cm}|p{3.5cm}|p{5cm}|}
\hline
\textbf{Campo} & \textbf{Tipo de dato} & \textbf{Restricciones} & \textbf{Propósito / Descripción} \\
\hline
id\_valoracion & SERIAL & PRIMARY KEY, NOT NULL & Identificador único de la valoración. \\
\hline
id\_receta & INT & NOT NULL, FOREIGN KEY & Referencia a la receta valorada. \\
\hline
id\_usuario & INT & NOT NULL, FOREIGN KEY & Referencia al usuario que realiza la valoración. \\
\hline
comentario & TEXT & NULL & Comentario opcional asociado a la valoración. \\
\hline
fecha\_comentario & TIMESTAMP & DEFAULT CURRENT\_DATE, CHECK & Fecha de la valoración; no puede ser posterior a la fecha actual. \\
\hline
valor & INT & CHECK (valor >= 1 AND valor <= 5) & Puntuación numérica de la receta. \\
\hline
\end{longtable}

\subsection{Tabla: CLASIFICACION}

Relaciona cada receta con uno o varios tipos de receta.

\begin{longtable}{|p{3.5cm}|p{2.5cm}|p{3.5cm}|p{5cm}|}
\hline
\textbf{Campo} & \textbf{Tipo de dato} & \textbf{Restricciones} & \textbf{Propósito / Descripción} \\
\hline
id\_receta & INT & NOT NULL, FOREIGN KEY & Referencia a la receta clasificada. \\
\hline
id\_tipo\_receta & INT & NOT NULL, FOREIGN KEY & Referencia al tipo de receta. \\
\hline
\end{longtable}

% ------------------------------------------------
% DOCUMENTACIÓN DE VISTAS
% ------------------------------------------------
\section{Documentación de Vistas}

\subsection{Vista: recetas\_mejor\_valoradas}

\textbf{Propósito:}

Presentar el top 5 de recetas con mayor promedio de valoración.

\textbf{Tablas relacionadas:}
\begin{itemize}
    \item recetas
    \item valoracion
    \item usuario
\end{itemize}

\textbf{Consulta SQL asociada:}
\begin{verbatim}
CREATE OR REPLACE VIEW recetas_mejor_valoradas AS
SELECT
    u.username            AS usuario,
    r.nombre_receta       AS receta,
    AVG(v.valor)          AS valoracion_promedio
FROM recetas r
JOIN valoracion v ON r.id_receta  = v.id_receta
JOIN usuario   u ON r.id_usuario  = u.id_usuario
GROUP BY r.id_receta, u.id_usuario
ORDER BY valoracion_promedio DESC
LIMIT 5;
\end{verbatim}

\textbf{Descripción funcional:}

Esta vista facilita el acceso a las recetas mejor valoradas sin tener que ejecutar manualmente los joins y agregaciones.

\subsection{Vista: reporte\_recetas}

\textbf{Propósito:}

Generar un reporte agrupado por tipo de receta con promedio de calificación y total de recetas.

\textbf{Tablas relacionadas:}
\begin{itemize}
    \item tipo\_receta
    \item clasificacion
    \item recetas
    \item valoracion
\end{itemize}

\textbf{Consulta SQL asociada:}
\begin{verbatim}
CREATE OR REPLACE VIEW reporte_recetas AS
SELECT
    tip.nombre_tipo       AS tipo_receta,
    AVG(v.valor)          AS promedio_valoracion,
    COUNT(r.id_receta)    AS total_recetas
FROM tipo_receta   tip
JOIN clasificacion c   ON tip.id_tipo_receta = c.id_tipo_receta
JOIN recetas       r   ON r.id_receta        = c.id_receta
JOIN valoracion    v   ON v.id_receta        = r.id_receta
GROUP BY tip.id_tipo_receta
ORDER BY total_recetas DESC;
\end{verbatim}

\textbf{Descripción funcional:}

Provee una vista de análisis para conocer el comportamiento de las recetas por categoría.

\subsection{Vista: usuario\_receta}

\textbf{Propósito:}

Mostrar los usuarios que poseen recetas sin exponer datos sensibles.

\textbf{Tablas relacionadas:}
\begin{itemize}
    \item usuario
    \item recetas
\end{itemize}

\textbf{Consulta SQL asociada:}
\begin{verbatim}
CREATE OR REPLACE VIEW usuario_receta AS
SELECT
    u.username
FROM usuario u
LEFT JOIN recetas r ON u.id_usuario = r.id_usuario;
\end{verbatim}

\textbf{Descripción funcional:}

Sirve para consultas de seguridad o informes donde sólo se requiere el nombre de usuario asociado a una receta.

\subsection{Vista: v\_recetas\_rapidas}

\textbf{Propósito:}

Filtrar recetas con tiempo de preparación menor o igual a 30 minutos.

\textbf{Tablas relacionadas:}
\begin{itemize}
    \item recetas
\end{itemize}

\textbf{Consulta SQL asociada:}
\begin{verbatim}
CREATE OR REPLACE VIEW v_recetas_rapidas AS
SELECT
    nombre_receta         AS receta,
    tiempo_preparacion    AS minutos,
    descripcion
FROM recetas
WHERE tiempo_preparacion <= 30
WITH CHECK OPTION;
\end{verbatim}

\textbf{Descripción funcional:}

Permite consultar rápidamente recetas de preparación corta y garantizar que los datos insertados o actualizados cumplan la condición definida.

% ------------------------------------------------
% CONCLUSIONES
% ------------------------------------------------
\section{Conclusiones}

La base de datos de RecipeBook fue diseñada siguiendo principios de integridad relacional y organización estructurada de la información, permitiendo un manejo eficiente de las entidades y relaciones involucradas en el sistema.

% ------------------------------------------------
% FIN
% ------------------------------------------------
\end{document}
